% kernel_loewner.tex -- round 494.
% PRIME.RDAGGER.KERNEL_LOEWNER.01.
% PROVED(c=2.1000e-3).  Companion: verify_kernel_loewner.py.
% Probe: experiments/tfpt-discovery/kernel_loewner_probe.py.
% Experiments-side compact-support Weil theorem.  NO RH CLAIM.
\documentclass[11pt]{article}

\usepackage[a4paper,margin=2.45cm]{geometry}
\usepackage{amsmath,amssymb,amsthm}
\usepackage{booktabs,array,microtype,xcolor}
\usepackage[T1]{fontenc}
\usepackage[colorlinks=true,linkcolor=blue!50!black,urlcolor=blue!50!black]{hyperref}

\newtheorem{theorem}{Theorem}
\newtheorem{lemma}{Lemma}
\newtheorem{remark}{Remark}
\newcommand{\QW}{Q_{\mathrm W}}
\newcommand{\proved}{\textcolor{green!40!black}{PROVED}}
\setlength{\emergencystretch}{2em}

\title{\bfseries Unconditional interval-certified compact-support
Weil positivity at $L=0.3$\\[2mm]
\large via the $x$-space kernel form\\[2mm]
\large PRIME.RDAGGER.KERNEL\_LOEWNER.01 --- round 494}
\author{Stefan Hamann}
\date{August 31, 2026}

\begin{document}
\maketitle

\begin{abstract}\noindent
The three-gate contract of the round-492 judgment is executed in its
prescribed order.  For every real
$h\in L^2(\mathbb R)$ supported in $[-0.3,0.3]$,
\[
  \QW(h)\ \ge\ 2.1\cdot10^{-3}\,\|h\|_2^2 .
\]
The proof is unconditional and interval-certified.  It uses the
classical $x$-space Weil kernel, positivity of every compressed
translation difference, a sign-safe cutoff, and one
$401$-dimensional Legendre finite section with an enclosed exact
Hilbert--Schmidt rest.  The Galerkin values of the uncut form are
calibration upper bounds only; none is used as an operator lower
bound.  This is an experiments-side compact-support theorem.  It is
not a claim for or against the Riemann Hypothesis.
\end{abstract}

\medskip\noindent
\textbf{Status.}\enspace
\texttt{PROVED(c=2.1000e-3)}.
No ledger row and no promotion of an RH claim.

\smallskip\noindent
Probe SPEC\_SHA:
\nolinkurl{2e6d71138dff454025641443a9c77f8018d36827dcf94cc35f1696dae9deda47}.

\begin{center}
\fbox{\parbox{0.95\linewidth}{%
\textbf{Claim boundary.}
The theorem is a fixed-support statement about the Weil quadratic
form at $L=0.3$.  It proves neither RH nor anti-RH and contains no
claim about zeros.  It does not assert that a finite Galerkin
eigenvalue is the operator eigenvalue.  The operator floor comes
from a Loewner minorant plus a complete finite-section remainder
bound.}}
\end{center}

\section{The theorem}

Let $\widetilde h(x)=h(-x)$ and $g=h*\widetilde h$.  For real $h$ put
\[
 \Pi(h)=2\left(
   \langle h,\cosh(\,\cdot\,/2)\rangle^2
  -\langle h,\sinh(\,\cdot\,/2)\rangle^2\right).
\]

\begin{theorem}[Unconditional interval-certified compact-support Weil
positivity at $L=0.3$ via the $x$-space kernel form]
\label{thm:kernel-loewner}
If $\operatorname{supp}h\subset[-0.3,0.3]$, then
\[
 \boxed{\quad \QW(h)\ge 2.1\cdot10^{-3}\,\|h\|_2^2.\quad}
\]
The same inequality holds for complex $h$ after applying the real
statement to its real and imaginary parts.
\end{theorem}

Since $2L=0.6<\log2$, the prime sum is empty.  Thus this theorem is
scramble-insensitive only vacuously; no source positions occur at
this support length.

\section{Gate G1: the sealed kernel identity}

\begin{lemma}[$x$-space Weil form]
\label{lem:xform}
For $2L<\log2$,
\begin{align}
 \QW(h)
 &=\int_0^{2L} k(x)\bigl(g(0)-g(x)\bigr)\,dx
       -c_Lg(0)+\Pi(h),                                      \label{eq:xform}\\
 k(x)&=\frac{e^{x/2}}{\sinh x},                              \label{eq:k}\\
 c_L&=\int_0^{2L}\left(k(x)-\frac1x\right)dx
       +\log(4L)+\gamma+\log\pi .                            \label{eq:cL}
\end{align}
At $L=0.3$,
\[
 c_L\in[2.19240491113,\,2.19240491114],
 \qquad c_L=2.1924049111319993023\ldots .
\]
\end{lemma}

\begin{proof}
For the defining archimedean multiplier, Gauss' integral for
$\psi(\tfrac14+it/2)$ and Plancherel give, after $v=2x$,
\[
 A(h)=\int_0^\infty
 \left(\frac{e^{-2x}}xg(0)-k(x)g(x)\right)dx
 -(\log\pi)g(0).
\]
The autocorrelation vanishes for $x>2L$.  Splitting there and using
\[
 \lim_{\varepsilon\downarrow0}\left(
 \int_\varepsilon^{2L}\frac{dx}{x}
 -\int_\varepsilon^\infty\frac{e^{-2x}}x\,dx\right)
 =\log(4L)+\gamma
\]
proves the archimedean part of \eqref{eq:xform}.
Furthermore
$\widehat g(i/2)+\widehat g(-i/2)
=2(\langle h,\cosh\rangle^2-\langle h,\sinh\rangle^2)$,
which fixes the pole sign, including the negative odd contribution.
The prime term is zero because $e^{2L}<2$.
\end{proof}

\paragraph{The three mandatory stop hits.}
The x-side regularization
\[
 \int_0^\infty\left[
 \frac{e^{-2x}-1}{x}-\left(k(x)-\frac1x\right)\right]dx-\log\pi
\]
gives
\[
 \sigma_A(0)=-5.3721834192256654.
\]
Its comparison with
$\psi(\tfrac14)-\log\pi$ differs by
$3.503\cdot10^{-26}$ after cutting the exponentially small tail at
$x=120$.  For
$h(x)=\cos(\pi x/(2L))$, direct insertion in \eqref{eq:xform} gives
\[
 \frac{\QW(h)}{\|h\|^2}=1.18025311091317\cdot10^{-2}
 \in[1.150,1.192]\cdot10^{-2}.
\]
The odd Legendre section gives
$2.2256359052\cdot10^{-1}$, consistent with the sealed
$0.211/0.221/0.226$ ladder.  Here
$\Pi_{\rm odd}=-2|\sinh(\,\cdot\,/2)\rangle
\langle\sinh(\,\cdot\,/2)|$; the opposite sign is not used.
All three stop conditions pass.

\section{Gate G2: Loewner truncation with boundary strips}

Let $E:L^2[-L,L]\to L^2(\mathbb R)$ be extension by zero and let
$\tau_x$ be the unitary translation
$(\tau_xf)(y)=f(y+x)$ on $L^2(\mathbb R)$.  Its interval compression
is $R_x=E^*\tau_xE$.  For $0\le x\le2L$,
\begin{align*}
 2\langle h,(I-\operatorname{Re}R_x)h\rangle
 &=\|Eh-\tau_xEh\|_{L^2(\mathbb R)}^2\\
 &=\int_{-L}^{L-x}|h(y)-h(y+x)|^2dy\\
 &\quad+\int_{-L}^{-L+x}|h(y)|^2dy
       +\int_{L-x}^{L}|h(y)|^2dy\ \ge0 .
\end{align*}
The last two integrals are the translation boundary strips.  They
would be lost if the norm were incorrectly taken only on
$[-L,L]$.  Consequently
\[
 I-\operatorname{Re}R_x\succeq0
 \quad\text{for every }x,
\]
and multiplication of the $x$-integral by any $0\le w\le1$ gives a
Loewner minorant.

The sealed cutoff is
\[
 w(x)=
 \begin{cases}
 0,&0\le x\le x_0,\\
 6t^5-15t^4+10t^3,&x_0<x<x_1,\quad
 t=(x-x_0)/(x_1-x_0),\\
 1,&x_1\le x\le2L,
 \end{cases}
 \qquad (x_0,x_1)=(0.01,0.03).
\]
This specific ramp is $C^2$ (not $C^\infty$); $C^2$ is sufficient
for the Hilbert--Schmidt certificate below.  Split 80-digit
quadrature gives
\[
 \kappa_w=\int_0^{0.6}w(x)k(x)\,dx
 =3.6980080402989774725\ldots
 \in[3.69800804029,3.69800804031].
\]

Define $T_w$ by the kernel
\[
 T_w(y,z)=\frac12w(|y-z|)k(|y-z|).
\]
Then
\[
 \QW\succeq \QW^{(w)}
  =(\kappa_w-c_L)I-T_w+\Pi.
\]
The omitted interval is therefore booked by using
\emph{$\kappa_w$}, never the uncut diagonal coefficient.  As an
audit diagnostic, the uncut $N=36$ Galerkin value
$7.5719610133\cdot10^{-3}$ falls to the truncated finite target
$4.2375222058\cdot10^{-3}$, a booked loss
$3.3344388074\cdot10^{-3}$, below half of the uncut diagnostic
margin.  This comparison is not an operator bound and is not used
in the proof.

\section{Gate G3: finite section and exact rest}

Let $P_n$ project onto the normalized Legendre polynomials of degrees
$0,\ldots,n-1$ on $[-L,L]$, with $n=401$.  Put
$S=T_w-\Pi$.  The enclosed finite-section value is
\[
 \lambda_{\max}(P_nSP_n)
 =1.501365606961153\ldots .
\]
Together with the outward surrogate and arithmetic radii,
\[
 \lambda_{\max}(P_nSP_n)\le1.5013673873 .
\]
Hence the direct interval lower bound before the block remainder is
\[
 \kappa_w^{\rm lo}-c_L^{\rm hi}
 -\lambda_{\max}^{\rm hi}(P_nSP_n)
 =4.2357418881\cdot10^{-3}.
\]

\subsection{Interval construction}

The square integral is changed to difference and overlap variables:
\[
 (T_w)_{ij}=\frac12\int_0^{2L}w(x)k(x)
 \int_{-L}^{L-x}
 \bigl(\phi_i(y+x)\phi_j(y)+\phi_i(y)\phi_j(y+x)\bigr)\,dy\,dx .
\]
Each of $[x_0,x_1]$ and $[x_1,2L]$ receives a degree-48 Chebyshev
surrogate.  Bernstein ellipses with respectively
$(\rho,M)=(2,15000)$ and $(1.5,220)$ give a full-kernel
Hilbert--Schmidt surrogate error
\[
 \varepsilon_{\rm sur}\le1.770301\cdot10^{-6}.
\]
For the surrogate the inner and outer Gauss--Legendre rules are
polynomial-exact.  Repeating with twelve extra nodes in both
directions changes the matrix by only
$2.0894\cdot10^{-11}$ in Frobenius norm and
$6.5745\cdot10^{-12}$ in operator norm; the sealed outward pads are
$5\cdot10^{-9}$ and $10^{-8}$.

The exact Hilbert--Schmidt identity is
\[
 r_n^2=
 \|T_w\|_{\rm HS}^2-\sum_{i,j<n}|(T_w)_{ij}|^2,\qquad
 \|T_w\|_{\rm HS}^2
 =\frac12\int_0^{2L}w(x)^2k(x)^2(2L-x)\,dx .
\]
Projection contraction transfers the surrogate enclosure to this
exact expression.  Its central value is
$6.7727140\cdot10^{-4}$ and the certified upper bound is
\[
 r_{401}\le7.0457565325\cdot10^{-4}
 <\frac13(4.2357418881\cdot10^{-3}).
\]
Thus the hard $n\le800$ NO-GO does not fire.

\subsection{The required threefold charge}

For any self-adjoint $S$,
\[
 \lambda_{\max}(S)\le\lambda_{\max}(P_nSP_n)
 +2\|P_nS(1-P_n)\|+\|(1-P_n)S(1-P_n)\|.
\]
For $T_w$, each of the three remainder norms is at most $r_n$.
The two off-diagonal blocks are therefore charged, not discarded:
\[
 r_n+2r_n=3r_n
 \le2.1137269598\cdot10^{-3}.
\]
The pole operator is rank two.  Taylor's remainder for
$\cosh(x/2)$ and $\sinh(x/2)$ after degree 400 bounds its complete
block remainder by $10^{-30}$.
Consequently
\[
 \QW^{(w)}
 \succeq
 \bigl(4.2357418881-2.1137269598\bigr)10^{-3}I
 =2.1220149285\cdot10^{-3}I
 \succeq2.1\cdot10^{-3}I .
\]
Together with $\QW\succeq\QW^{(w)}$, this proves
Theorem~\ref{thm:kernel-loewner}.

\section{Budget and Galerkin audit}

The requested human-readable audit decomposition is
\[
\underbrace{7.5719610133}_{\text{uncut Galerkin diagnostic}}
-\underbrace{3.3344388074}_{\text{booked Loewner loss}}
-\underbrace{2.1137269598}_{\text{finite-section block rest}}
=2.1237952461
\quad(\times10^{-3}).
\]
It explains the measured budget but is not the logical lower-bound
chain, because the first term is a Galerkin upper bound.  The
load-bearing operator chain is the direct interval balance in the
preceding section and yields $2.1220149285\cdot10^{-3}$.

\begin{center}
\small
\begin{tabular}{@{}rrrrrrr@{}}
\toprule
$N$ & 6 & 10 & 14 & 18 & 24 & 36\\
\midrule
$\lambda_{\min}^{\rm even}\cdot10^3$
&7.596752&7.580831&7.575943&7.573987&7.572755&7.571961\\
\bottomrule
\end{tabular}
\end{center}
The sequence is monotone from above; its spread from $N=6$ to
$N=36$ is $2.4791\cdot10^{-5}$.  Every entry remains explicitly
typed as a finite-section upper bound.

\section{Kill and pad audit}

\begin{center}
\small
\begin{tabular}{@{}p{0.30\linewidth}p{0.59\linewidth}@{}}
\toprule
Consolidated family & R494 disposition\\
\midrule
F-I high-band majorants & absent: the proof uses positive translation
differences and one compact kernel\\
F-II support-mismatched bases & absent: Legendre functions live on the
support interval and only certify a finite section\\
F-III unresolved discrete modes & absent: the exact complementary
block is enclosed by $r_{401}$\\
F-IV quadrature/truncation pads & absent: the cutoff loss is booked and
the surrogate has an outward Bernstein enclosure\\
F-V wrong comparison objects & absent: no shifted comparison,
Schur pad, or asymmetric product enters\\
F-VI false operator premises & absent: no inverse or coercivity premise
is used\\
\bottomrule
\end{tabular}
\end{center}

\begin{center}
\small
\begin{tabular}{@{}p{0.06\linewidth}p{0.83\linewidth}@{}}
\toprule
N1 & blocked: Galerkin numbers are labelled upper bounds only\\
N2 & blocked: $[0,x_0]$ is booked through $\kappa_w$ and the loss line\\
N3 & blocked: the rest is $r+2r=3r$, including both off-block terms\\
N4 & blocked: after the mandatory G1 identity calibration, the proof
stays entirely in the $x$-space translation/kernel representation\\
\bottomrule
\end{tabular}
\end{center}

\section{Scaling beyond the prime-free rung}

When $2L\ge\log2$, the prime term is no longer zero.  It enters the
same representation as the explicit bounded translation sum
\[
 -\sum_{n\le e^{2L}}\frac{\Lambda(n)}{\sqrt n}
   \bigl(R_{\log n}+R_{-\log n}\bigr).
\]
At $L=0.8$ there are exactly three terms, $n=2,3,4$.  Each compressed
translation has norm at most one, so this is a finite bounded
perturbation.  The Loewner monotonicity of the positive
archimedean translation-difference integral is unchanged.  Unlike
the present $L=0.3$ rung, the literal nodes $\log n$ make the
scramble gate live there.

\section{What is not claimed}

No Galerkin eigenvalue is promoted to an operator value.  No result
for every support length is asserted.  No zero statistic or zero
oracle is used.  The theorem is an unconditional fixed-support
Weil-positivity statement in the classical Yoshida--Bombieri zone,
not an RH claim and not an anti-RH claim.

\end{document}
